\section{Common disks for positive proportions of branches}
\label{sec:common-deformation}
The direction theorem concerns derivatives at the initial zeros. We
now continue fixed positive proportions of the actual branches over a
single parameter disk. A stronger derivative estimate also enlarges
the disk on which these branches remain simple and disjoint.

\begin{theorem}\label{thm:common-deformation}
There are absolute constants $c_0,c_1>0$ with the following property.
For every fixed primitive character $\chi$ and all sufficiently large
$T$ depending on $\chi$, each strict normal-velocity sign has at least
$10^{-24}N_{\chi,T}$ distinct simple centers $\rho$ in $(T,2T)$.
Their branches in the family $F_\chi(u,s)$ of
Section~\ref{sec:signed-density} are holomorphic, simple and pairwise
disjoint throughout
\[
                   |u|\le c_1\frac{\sqrt{\log\log T}}T.
\]
For every positive real $u\le c_0/T$, the same two sets of branches
lie on their respective strict sides of the critical line. All these
branches have ordinates in $(T,2T)$ throughout the complex disk.
\end{theorem}

\subsection{Weighted curvature}
Let $B=B_{\chi,T}$, $D_\rho=L_\chi'(\rho)$, and put
\[
 R=\frac1{1000L},\qquad
 M_\rho=\max_{|z|\le R/2}|L_\chi''(\rho+z)|.
\]
There are absolute constants $C_Q,C_M$ such that, eventually for each
fixed character,
\begin{equation}\label{eq:weighted-curvature-new}
 \sum|B(\rho)|^4|L_\chi(\rho+z)|^2\le C_QN
       \quad(|z|\le R),\qquad
 \sum|B(\rho)|^4M_\rho^2\le C_MNL^4.
\end{equation}
To prove the first assertion use the undifferentiated shifted moment
\eqref{eq:dirichlet-discrete}, with its square-mollifier coefficients.
Take outer physical shift circles of radii $2/(1000L)$ and
$3/(1000L)$, and a smaller auxiliary circle. The separate pole
denominators stay apart on these circles. The estimates in
\eqref{eq:square-euler-F}--\eqref{eq:square-euler-G} give $O(TL)$:
the first coefficient product is $O(L^{-1})$, its explicit pole is
$O(L)$ and the gamma logarithm is $O(L)$; the second product is
$O(1)$ and its auxiliary derivative is $O(L)$.
The Mellin tails were bounded in the proof of
Proposition~\ref{prop:square-moments}. The conductor residues cancel,
so the limiting constants are absolute; the lower-order errors have a
character-dependent starting height. The combined finite shifted sum
is holomorphic. The maximum principle extends the bound to both
interior disks. Set the shifts to $z,\bar z$; reflection on the
critical line makes this sum precisely the positive one in
\eqref{eq:weighted-curvature-new}. Since $N\sim TL/(2\pi)$ this
proves the assertion. Cauchy's formula gives
\[
 M_\rho^2\le256R^{-4}\frac1{2\pi}
                  \int_0^{2\pi}|L_\chi(\rho+Re^{it})|^2dt.
\]
Multiplying by $|B|^4$, summing and using the first assertion proves
the second.

\subsection{A logarithm sampled at separated centers}
Fix $0<\nu_0<1/100$, $x=T^{\nu_0}$ and $\eta>0$. Let
\[
 w_x(n)=\begin{cases}
 1,&n\le x,\\
 \log(x^2/n)/\log x,&x<n\le x^2,\\
 0,&n>x^2,
 \end{cases}
 \qquad
 P^*_{\chi,T}(t)=\sum_{2\le n\le x^2}
 \frac{\Lambda(n)w_x(n)\chi(n)}{n^{1/2+\eta/L+it}\log n}.
\]
The needed logarithmic estimate, with every ambient multiplicity, is
\begin{equation}\label{eq:log-kernel-new}
 \left|\log|L_\chi(1/2+\eta/L+it)|-\Re P^*_{\chi,T}(t)\right|
 \le C_\eta E_T(t)+1,\qquad
 E_T(t)=\sum_{\gamma'}\frac1{1+L^2(t-\gamma')^2}.
\end{equation}
The sum includes all positive, negative and zero ordinates of the
same character. One can take
$C_\eta=\max(4/(\nu_0\eta),3/\nu_0^2)$.

For completeness, integrate
\[
 \frac1{2\pi i\log x}\int_{(2)}
 -\frac{L_\chi'}{L_\chi}(s+w)
                    \frac{x^{2w}-x^w}{w^2}\,dw,
       \qquad \Re s>1/2.
\]
Perron's elementary $w^{-2}$ integral gives the smoothed prime sum.
Shift to the left along lines midway between the trivial zeros. The
functional equation bounds the remaining integral by
$O_\chi(x^{-A}(\log(A+|t|+3)+1)/A)$ on a line at distance $A$
to the left. Good horizontal edges tend to zero by the local zero
count and the squared denominator. The zero residues converge
absolutely, giving
\begin{align*}
 -\frac{L_\chi'}{L_\chi}(s)
 ={}&\sum_n\Lambda(n)w_x(n)\chi(n)n^{-s}
 +\frac1{\log x}\sum_{\omega\ {\mathrm{zero}}}
          \frac{x^{2(\omega-s)}-x^{\omega-s}}{(\omega-s)^2}\\
 &-\frac{\mathbf1_{q=1}}{\log x}
          \frac{x^{2(1-s)}-x^{1-s}}{(1-s)^2}.
\end{align*}
Integrate horizontally from $\Re s=1/2+\eta/L$ to infinity.
The Euler logarithm determines the branch; the zero-location input
makes the ray zero-free. For a nontrivial zero put
$v=L(\sigma-1/2)$ and $y=L(t-\gamma')$. Its integrated absolute
contribution is at most
\[
 \frac1{\nu_0}\int_\eta^\infty
       \frac{e^{-\nu_0v}+e^{-2\nu_0v}}{v^2+y^2}\,dv
                \le\frac{C_\eta}{1+y^2}.
\]
For $|y|\le1$ use the integral of $2v^{-2}$; for $|y|\ge1$
integrate the exponentials divided by $y^2$. The trivial-zero sum
is $O(x^{-1/2}/(T^2(\log x)^2))$ and the zeta pole costs at most
$2x/(T^2(\log x)^2)$. Both together are eventually below one.
Taking real parts proves \eqref{eq:log-kernel-new} and also
justifies the interchanges by nonnegative domination.

We will use two elementary sampling consequences. First,
\begin{equation}\label{eq:prime-comparison-new}
 \sum_{T<\gamma\le2T}
    |P^*_{\chi,T}(\gamma)-P_T(\rho)|^2=o(NV),
\end{equation}
where $P_T$ is the marked prime polynomial of
Section~\ref{sec:gaussian-moments}. To check this, split the difference
into omitted primes, the change $p^{-\eta/L}-1$ at included primes,
and higher prime powers. Their squared coefficient norms with weight
$1/n$ are, respectively,
$O_\chi(1+(\log V)^2)$, $O(d_T^2V)$, and $O(1)$.
For a prime-supported piece, two distinct primes have no integer
prime-power ratio, so its off-diagonal main terms in
\eqref{eq:rational-landau-new} vanish. For the higher prime powers
the coefficients are nonnegative; the integer-ratio main term in the
expanded square is
\[
             -\frac T\pi\frac{a_na_m\Lambda(m/n)}m\le0
             \quad(m>n).
\]
The character factors cancel when $(m,q)=1$, and otherwise the
term is zero. This proves the sign also for complex characters.
The absolute errors, with $Y=x^2$, are bounded by
\[
 L^2\sum_{n<m\le Y}\frac1n\ll YL^3,\quad
 \sum_{n<m\le Y}\frac{\log(2m/n)}
                  {n\operatorname{dist}_*(m/n,\mathbb Z)}
                       \ll Y^2L,\quad
 L\sum_{n<m\le Y}\frac1{m\log(m/n)}\ll YL^2.
\]
Use $\operatorname{dist}_*\ge1/n$ and
$\log(m/n)\ge(m-n)/m$. They are $o(N)$; the endpoint strips cost
$O_\chi(YL)=o(N)$. Apply the square inequality to the three pieces.
Since $(\log V)^2=o(V)$, this proves
\eqref{eq:prime-comparison-new}.

Second, for every set $\mathcal S$ of distinct centers in $(T,2T)$
with mutual separation at least $b/L$,
\begin{equation}\label{eq:reversed-kernel-new}
                     \sum_{\gamma\in\mathcal S}E_T(\gamma)
                          \le C(b)N.
\end{equation}
For each real $u$, separation gives
$\sum_{\gamma\in\mathcal S}(1+L^2(\gamma-u)^2)^{-1}\le C(b)$
by comparison with the convergent lattice sum.
Reverse the sum over ambient zeros in $[T/2,3T]$, whose multiplicity
is $O(N)$. The remaining zeros with $|\gamma'|\le4T$ contribute
$O(\#\mathcal S/(TL))$, since their distance is at least $T/2$.
On the shells $2^jT<|\gamma'|\le2^{j+1}T$, $j\ge2$, the count
is $O(2^jT(L+j))$ and the distance at least $2^{j-1}T$.
Their contribution is bounded by
\[
 \frac{C\#\mathcal S}{TL^2}\sum_{j\ge2}2^{-j}(L+j)
                 \ll\frac{\#\mathcal S}{TL}.
\]
These counts have absolute eventual constants for each fixed
character; negative ordinates are counted using its conjugate.
Tonelli's theorem applies throughout. This proves
\eqref{eq:reversed-kernel-new} without removing any complementary
multiple zero.

\subsection{Localization for every fixed power of the conductor}

The logarithmic comparison itself admits every fixed polynomial
conductor range. We record its quantifiers separately from the
moment estimates and the selection of signed sets. All definitions of
$P_T,\sigma_T,X_\rho,R,M_\rho$ remain unchanged.

\begin{proposition}\label{prop:uniform-separated-log-localization}
Fix $A,h,K>0$, a bounded interval $J$, and $\epsilon>0$.
Uniformly over primitive characters of conductor $q\le T^A$,
including conductor one, and all sets $\mathcal S_T$ of distinct
simple centers in $(T,2T)$ satisfying
\[
 |\gamma-\gamma'|\ge h/L\quad(\rho\ne\rho'),\qquad
 M_\rho\le KL|D_\rho|\quad(\rho\in\mathcal S_T),
\]
one has
\begin{align}
 &\#\left\{\rho\in\mathcal S_T:X_\rho\in J,
 \left|\frac{\log|D_\rho|-\log L}{s_T}-X_\rho\right|>\epsilon
 \right\}\nonumber\\
 &\hspace{25mm}\ll_{A,h,K,J,\epsilon}
 N_{\chi,T}\left(V^{-1/4}+\frac{(\log V)^2}{V}\right),
 \qquad s_T=\sqrt{V/2}.\label{eq:polynomial-conductor-localization}
\end{align}
The starting height is uniform in the character and in the choice of
$\mathcal S_T$. Every ambient zero count includes multiplicity.
\end{proposition}

\begin{proof}
Put $Q=\log(qT)$, so $L\le Q\le(1+A)L$, and set
\[
 W=V^{1/4},\qquad y=T^{1/W},\qquad a=\frac1{\log y}=\frac WL,
 \qquad c=\min(1/4000,1/(4K)),\qquad v=c/L.
\]
Write
\[
 F_u(t)=\sum_{\gamma'}\frac{u}{u^2+(t-\gamma')^2}.
\]
The centered Hadamard identity proved in
Section~\ref{sec:dirichlet-moments}, with the gamma factor removed,
gives uniformly for $0<u\le1$ and $T\le t\le2T$
\begin{equation}\label{eq:localization-real-log-derivative}
 \Re\frac{L_\chi'}{L_\chi}(1/2+u+it)
     =F_u(t)-\frac12\log\frac{qt}{2\pi}+O(T^{-1}).
\end{equation}
For complex characters the exact reflection is
$\Lambda_\chi(s)=\varepsilon_\chi
\overline{\Lambda_\chi(1-\bar s)}$; thus the centered linear
coefficient has zero real part. For conductor one use the entire
completion with its factors $s(s-1)$. This proves the displayed
identity for the same character without pairing its positive and
negative ordinates.

At $u=1$ the Euler logarithmic derivative is absolutely bounded,
so $F_1(t)=\tfrac12\log(q(|t|+2))+O(1)$ for all real $t$.
Consequently each unit interval at height $t$ contains
$O(\log(q(|t|+2)))$ occurrences. These bounds also give
\begin{equation}\label{eq:localization-occurrence-count}
 N:=N_{\chi,T}\asymp TQ,\qquad
 \#\{\gamma':T/2\le\gamma'\le3T\}\ll N.
\end{equation}
Here is a direct lower bound for the first assertion. For
$5T/4\le t\le7T/4$, dyadic summation of the local count gives
\[
 \sum_{\gamma'\notin(T,2T]}
       \frac1{1+(t-\gamma')^2}\ll Q/T.
\]
Indeed a distance shell of size $2^jT$ contains
$O(2^jT(Q+j+1))$ occurrences. Integrating $F_1$ over this middle
interval gives $TQ/4+O(T)$; the outside occurrences contribute
$O(Q)$ and each inside occurrence contributes at most $\pi$.
Thus $\pi N\ge TQ/4-O(T+Q)$. The upper bound and the second
assertion follow by covering with unit intervals. This proves
\eqref{eq:localization-occurrence-count} uniformly in the present
range.

List the ordinates in $\mathcal S_T$ as $t_j$. For an even
nonnegative function decreasing on the positive ray, separation
$\delta=h/L$ gives
\[
 \sum_j f(t_j-b)\le2f(0)+\delta^{-1}\int_{\mathbb R}f(t)\,dt.
\]
Retain the nearest point on each side of $b$, and bound each
remaining value by the integral over the preceding interval of length
$\delta$. Apply this to $f(t)=u/(u^2+t^2)$ and sum over the
occurrences in $[T/2,3T]$. The complementary occurrences have
$\sum(t_j-\gamma')^{-2}\ll Q/T$ by the same dyadic argument.
Since $\#\mathcal S_T\le N$, we obtain
\begin{equation}\label{eq:localization-summed-kernel}
 \sum_jF_u(t_j)\ll
 N\left(u^{-1}+L/h\right)+NuQ/T\qquad(0<u\le1).
\end{equation}

Define the auxiliary polynomial
\[
 \mathcal Q_y(t)=\sum_{2\le n\le y}
 \frac{\Lambda(n)\chi(n)}{n^{1/2+a+it}\log n}
                       \frac{\log(y/n)}{\log y}.
\]
The smoothed explicit formula follows the method of
\cite[Lemma~1 and its integrated form, pp.~985--987]{soundararajan-moments2009}
and \cite[equations~(2.2)--(2.6), pp.~111--112]{sono-dirichlet2013};
we retain an absolute error in both directions.
The $w^{-2}$ contour used in the proof of
\eqref{eq:uniform-character-growth} gives the following integrated
identity, with $s_a=1/2+a+it$:
\begin{align}
 \log|L_\chi(s_a)|
 ={}&\Re\mathcal Q_y(t)
       -a\Re\frac{L_\chi'}{L_\chi}(s_a)+\Re Z_y(t)+R_y(t),
       \label{eq:localization-smoothed-log}\\
 Z_y(t)={}&\frac1{\log y}\sum_{\gamma'}\int_a^\infty
 \frac{y^{-u+i(\gamma'-t)}}{(-u+i(\gamma'-t))^2}\,du.
 \nonumber
\end{align}
To check every term, integrate
$-(L_\chi'/L_\chi)(s+w)y^w/w^2$ on a right vertical line.
The residue at zero is
$-(L_\chi'/L_\chi)(s)\log y-(L_\chi'/L_\chi)'(s)$;
a zero of multiplicity $m$ contributes
$-m y^{\rho-s}/(\rho-s)^2$, and the zeta pole contributes
$y^{1-s}/(1-s)^2$. Shift left as in that earlier contour proof,
divide by $\log y$, and integrate horizontally from $s_a$ to
infinity. The derivative term gives the second term in
\eqref{eq:localization-smoothed-log}. The trivial-zero and pole
integrals give
\[
 |R_y(t)|\ll
 \frac{y^{-1/2-a}+\mathbf1_{q=1}y^{1/2-a}}
                 {T^2(\log y)^2}\ll T^{-1}.
\]
This follows by summing the parity progression of trivial zeros,
using denominators at least $T^2$, and integrating the exponential
numerators. All constants are absolute. The nontrivial-zero integral
satisfies
\[
 |Z_y(t)|\le
 \frac{y^{-a}}{(\log y)^2}
       \sum_{\gamma'}\frac1{a^2+(t-\gamma')^2}
             =e^{-1}aF_a(t).
\]
This absolute majorant justifies the integration and interchange of
the zero sum. The zero-free horizontal ray is supplied by the
admitted zero-location theorem.
Equations~\eqref{eq:localization-real-log-derivative} and
\eqref{eq:localization-summed-kernel} now give
\begin{equation}\label{eq:localization-auxiliary-log-mean}
 \sum_j\left|\log|L_\chi(1/2+a+it_j)|
                      -\Re\mathcal Q_y(t_j)\right|
                      \ll_{A,h}NW.
\end{equation}

The bridge to $v$ uses the exact integral of
\eqref{eq:localization-real-log-derivative}. In particular,
\[
 \int_v^aF_u(t)\,du
 =\frac12\sum_{\gamma'}
       \log\frac{a^2+(t-\gamma')^2}{v^2+(t-\gamma')^2}.
\]
The integrands are nonnegative. Integrating
\eqref{eq:localization-summed-kernel} therefore shows
\begin{align*}
 &\sum_j\left|\log|L_\chi(1/2+a+it_j)|
                  -\log|L_\chi(1/2+v+it_j)|\right|\\
 &\quad\ll N\{\log(a/v)+L(a-v)/h+(a^2-v^2)Q/T+aQ\}
            \ll_{A,h,K}NW.
\end{align*}
Only the displacement $v$ is used for Taylor's formula. It lies
inside $R/2$, and the stated curvature condition gives
\[
 L_\chi(\rho+v)=vD_\rho(1+e_\rho),\qquad
 |e_\rho|\le vM_\rho/(2|D_\rho|)\le Kc/2\le1/8.
\]
Thus, with
$U_\rho=\log|D_\rho|-\log L-\Re\mathcal Q_y(\gamma)$,
\begin{equation}\label{eq:localization-derivative-log-mean}
             \sum_{\mathcal S_T}|U_\rho|\ll_{A,h,K}NW.
\end{equation}
The bounded extra logarithmic factors are
$-\log c-\log|1+e_\rho|$.

We next compare the literal coefficients of $\mathcal Q_y$ and
$P_T$. Eventually $z_T<x_T<y$. For a prime $p\le y$, put
$r_p=\log p/\log y$ and $w_p=e^{-r_p}(1-r_p)$. The difference
coefficient is exactly
\[
 \frac{\chi(p)}{\sqrt p}
       (w_p-\mathbf1_{z_T<p\le x_T}).
\]
Use $0\le w_p\le1$, $|w_p-1|\le2r_p$, and the
reciprocal-prime estimates of Section~\ref{sec:gaussian-moments}.
Partial summation also gives
$\sum_{p\le u}(\log p)^2/p\ll(\log u)^2$.
The full coefficient square sum is consequently at most
\begin{align*}
 &\sum_{p\le z_T}\frac1p
 +\frac4{(\log y)^2}\sum_{p\le x_T}\frac{(\log p)^2}p
 +\sum_{x_T<p\le y}\frac1p
 +\sum_p\sum_{j\ge2}\frac1{j^2p^j}\\
 &\quad\ll 1+\log\log z_T+(d_TW)^2+\log(1/(d_TW))
       \ll(\log V)^2.
\end{align*}
The prime-power coefficient is bounded by $1/(jp^{j/2})$.
These bounds retain $\chi(p^j)=\chi(p)^j$ and use only
$|\chi(n)|\le1$, uniformly in the conductor.

Here a direct separated sampling inequality suffices. For
$H(t)=\sum_{n\le y}b_nn^{-it}$ and $E=\sum|b_n|^2$,
expansion on any interval of length at most $2T$ gives
\[
 \int|H|^2\ll(T+y\log(2y))E,\qquad
 \int|H'|^2\ll(T+y\log(2y))(\log y)^2E.
\]
For the off-diagonal terms use
$2|b_nb_m|\le|b_n|^2+|b_m|^2$ and
$\sum_{m\ne n}|\log(m/n)|^{-1}\ll y\log(2y)$,
which follows from
$|\log(m/n)|\ge|m-n|/\max(m,n)$.
The intervals $I_j=[t_j-\delta/2,t_j+\delta/2]$ have disjoint
interiors. Averaging the fundamental theorem of calculus on each
one gives
\[
 |H(t_j)|^2\le\delta^{-1}\int_{I_j}|H|^2
                             +2\int_{I_j}|HH'|.
\]
Sum and apply Cauchy's inequality. Since $y\log(2y)=o(T)$ and
$\log y\le L$, the result is
$\sum_j|H(t_j)|^2\ll_h TLE\ll_h NE$.
Applying it to the coefficient difference proves
\begin{equation}\label{eq:localization-polynomial-square-mean}
 \sum_{\rho\in\mathcal S_T}
 |\mathcal Q_y(\gamma)-P_T(\rho)|^2\ll_hN(\log V)^2.
\end{equation}

Finally, the variance has its stated normalization uniformly in
this larger conductor range. Set $H_q=\max(3,\log(q+3))$.
The prime divisors larger than $H_q$ contribute at most
$\log q/(H_q\log H_q)$ to $\sum_{p\mid q}p^{-1}$;
the others contribute at most $\log\log H_q+O(1)$.
Thus the omitted primes contribute $O_A(1+\log V)$, and
\[
 2\sigma_T^2
 =V-(\log V)^2-2\log\log V+O_A(1+\log V),\qquad
 \frac{\sigma_T}{s_T}=1+O_A((\log V)^2/V).
\]
Put $H_\rho=\mathcal Q_y(\gamma)-P_T(\rho)$. There is the
exact equality
\[
 \frac{\log|D_\rho|-\log L}{s_T}-X_\rho
 =\frac{U_\rho+\Re H_\rho}{s_T}
              +X_\rho(\sigma_T/s_T-1).
\]
On $X_\rho\in J$ the last term is at most $\epsilon/2$ eventually,
uniformly. Each exception therefore has
$|U_\rho|>\epsilon s_T/4$ or $|H_\rho|>\epsilon s_T/4$.
Markov's inequality, \eqref{eq:localization-derivative-log-mean},
and \eqref{eq:localization-polynomial-square-mean} bound their number by
\[
 \frac4{\epsilon s_T}\sum|U_\rho|
 +\frac{16}{\epsilon^2s_T^2}\sum|H_\rho|^2
 \ll_{A,h,K,\epsilon}
 N\left(V^{-1/4}+\frac{(\log V)^2}{V}\right).
\]
This proves the proposition with all the asserted uniformities.
\end{proof}

\subsection{Selection with a large derivative}
\begin{lemma}\label{lem:large-derivative-selection}
There are absolute $K_0\ge1$, $h>0$, $\eta>0$, and two sets of
simple centers, each containing at least $10^{-24}N$ points, such that
\begin{gather*}
 T+1<\Im\rho<2T-1,\qquad
 M_\rho\le K_0L|D_\rho|,\qquad
 |D_\rho|\ge\frac L{2\eta}e^{\sigma_T/2-1},\\
 \frac{|\Re v_\rho|}{|v_\rho|}\ge10^{-12},\qquad
 v_\rho=\frac{\rho L_\chi(\rho+1)}{D_\rho}.
\end{gather*}
The two sets have opposite strict signs of $\Re v_\rho$, and all
their centers together have mutual separation at least $4h/L$.
\end{lemma}

\begin{proof}
Restrict the five interval limits to $I=[1,2]$. Its Gaussian mass
$p$ exceeds $1/24$, since $p\ge e^{-2}/\sqrt{2\pi}$.
The first two mollifier masses become $apNL+o(NL)$ and
$bpNL^2+o(NL^2)$; the phase mass has
$\sum_I\cos^2\psi=(p/2+o(1))N$.
With $\eta_0=1/50$, the bad phase set has size at most
$(p/(2(1-\eta_0^2))+o(1))N$. The same exact numerical margin
as in the direction theorem gives
\[
 \sum_{I,\,|\cos\psi|\ge\eta_0}|A_\rho|
     \ge\left(p\left(a_*-\sqrt{b_*/(2(1-\eta_0^2))}\right)
                                                   -o(1)\right)NL
     >(p/4000-o(1))NL.
\]
The Euler bounds $1/3<|L_\chi(3/2+i\gamma)|<3$ and the exact
phase identity give total absolute $C_\rho$ mass greater than
$(p/600000-o(1))NL$. Its signed sum is $o(NL)$ by the fifth interval
limit. Each strict sign therefore has mass exceeding $NL/40000000$
eventually. Also $\sum|A_\rho|^2\le NL^2$ and
$|C_\rho|\le6|A_\rho|$ eventually.

Put $\kappa=10^{-9}$ and $\lambda=10^{-12}$.
Discard $|\cos\psi|<\lambda$, which costs at most $6\lambda NL$.
The unit endpoint strips have $O_\chi(L)$ occurrences and therefore
cost $o(NL)$ by Cauchy. Choose
$K_0\ge100\sqrt{C_M}/\kappa$. On the set
$M_\rho>K_0L|D_\rho|$, the absolute signed mass is at most
\[
 \frac6{K_0L}\sum |B(\rho)|^2M_\rho
 \le\frac{6\sqrt{C_M}}{K_0}NL<\kappa NL/8.
\]
Multiple centers carry zero mass and are removed. Choose
$h\le\min(1/16000,1/(100K_0))$.
Taylor's formula at any remaining center shows there is no other
such center within $4h/L$: otherwise
$|D_\rho|\le2K_0h|D_\rho|$, a contradiction. The segment lies in
the defining curvature disk.

Fix $\eta\le\min(1/4000,1/(4K_0))$ and apply
\eqref{eq:log-kernel-new}. By \eqref{eq:prime-comparison-new},
$|P^*_{\chi,T}-P_T|\le\sigma_T/4$ fails at only $o(N)$ occurrences.
By \eqref{eq:reversed-kernel-new},
$C_\eta E_T\le\sigma_T/4$ fails at only $O(N/\sigma_T)=o(N)$
remaining centers. Their total absolute signed mass is $o(NL)$ by
Cauchy. Since $X_\rho\in[1,2]$ at every center under consideration,
\[
          \log|L_\chi(\rho+\eta/L)|\ge\sigma_T/2-1.
\]
Taylor's formula gives
$L_\chi(\rho+\eta/L)=(\eta/L)D_\rho(1+e_\rho)$ with
$|e_\rho|\le K_0\eta/2\le1/8$. Hence the displayed exponential
lower bound for $|D_\rho|$ holds.
Each sign still has absolute mass at least $\kappa NL$.
Cauchy and $\sum C_\rho^2\le36NL^2$ give at least
$(\kappa^2/36)N>10^{-24}N$ distinct centers of each sign.
All other asserted conditions were retained in the deletions.
\end{proof}

\subsection{The two continuation scales}
\begin{proof}[Proof of Theorem~\ref{thm:common-deformation}]
The exact denominator-shift series is
\[
 F_\chi(u,s)=L_\chi(s)+
        \sum_{k\ge1}\frac{(-u)^k(s)_k}{k!}L_\chi(s+k).
\]
It follows first from the binomial series in $\Re s>1$, then by
analytic continuation; its terms with $k\ge1$ are locally normally
convergent near the high centers for $|u|<1$.
Put $G(s)=sL_\chi(s+1)$. On fixed disks of radius less than $1/100$
around those centers, absolute Dirichlet series give
\[
 c_GT\le|G(\rho)|\le C_GT,\qquad |G'(s)|\le C_GT.
\]
All constants are absolute. The lower bound uses the Euler product
at real part $3/2$. For $v=|s|\le3T$,
$|(s)_k|\le(v)_k$ and
$\sum_{k\ge0}(v)_k|u|^k/k!=(1-|u|)^{-v}$.
Thus, when $|u|\le1/10$, the whole perturbation satisfies
\begin{equation}\label{eq:shift-exponential-bound}
 |F_\chi(u,s)-L_\chi(s)|+
       |\partial_sF_\chi(u,s)-L_\chi'(s)|
       \le C T|u|e^{C T|u|}.
\end{equation}
Cauchy's formula on a slightly larger fixed $s$ disk gives the
derivative bound. The terms of order at least two separately give
\begin{equation}\label{eq:shift-quadratic-bound}
 F_\chi(u,s)=L_\chi(s)-uG(s)+R(s,u),\qquad
 |R|+|R_s|\le C_RT^2|u|^2
       \quad(T|u|\le1/10).
\end{equation}

Keep the sets in Lemma~\ref{lem:large-derivative-selection} and the
pairwise disjoint center disks $|s-\rho|<r=h/L$.
On their boundaries,
\[
 |L_\chi(s)-D_\rho(s-\rho)|
       \le\tfrac12M_\rho r^2\le(K_0h/2)|D_\rho|r.
\]
Choose an absolute $d>0$ with $Cd<1/4$, where $C$ is from
\eqref{eq:shift-exponential-bound}. Uniformly for $|u|\le d\sigma_T/T$,
the ratio of the perturbation to $|D_\rho|r$ is at most
\[
 C'\sigma_T\exp\{(Cd-1/2)\sigma_T\}\longrightarrow0.
\]
Rouch\'e's theorem gives exactly one zero in each disk for every
such complex parameter. Moreover throughout each disk,
\[
 |\partial_sF_\chi(u,s)-D_\rho|
 \le K_0h|D_\rho|+CT|u|e^{CT|u|}<|D_\rho|/2
\]
eventually. The zero is simple. Its local implicit sections patch by
uniqueness to a holomorphic section throughout the parameter disk.
Since $\sigma_T\sim\sqrt{\log\log T/2}$, an absolute $c_1>0$
gives the claimed disk. The fixed center disks are disjoint and lie
between the two discarded ordinate strips, proving the simultaneous
distinctness and height assertions on that whole disk.

We establish the uniform normal direction on the stated positive
real interval. Let $s_\rho(u)=\rho+z(u)$. For $T|u|\le c_0$ with
$c_0$ small and absolute, the zero equation and
\eqref{eq:shift-quadratic-bound} first give
$|z(u)|\le C|u|T/|D_\rho|$. Subtract the linear displacement
$uv_\rho=uG(\rho)/D_\rho$ from that equation. For $u\ne0$,
\[
 \left|\frac{z(u)}{uv_\rho}-1\right|
 \le C'\left(
       \frac{M_\rho|u|T}{|D_\rho|^2}
       +\frac{|u|T}{|D_\rho|}+T|u|\right)
 \le C'((K_0+1)c_0+c_0/L),
\]
where we used the weaker consequence $|D_\rho|\ge L$.
Choose the same $c_0$ small enough that the final bound is less than
$10^{-12}/2$. For positive real $u$ the margin
$|\Re v_\rho|\ge10^{-12}|v_\rho|$ then forces $\Re z(u)$ to have
the prescribed strict sign. This establishes the second scale and
completes the proof.
\end{proof}

Every selected zero defines a degree-one analytic cover over the
common complex disk: the derivative in the zero coordinate is
nonzero everywhere and the unique section has just been constructed.
The two normal orientations are the signs of the real differential
of this same section on the specified positive parameter ray.
